\documentclass{article}
\usepackage[utf8]{inputenc}
\usepackage[brazil]{babel}
\usepackage{geometry}
\usepackage{titlesec}
\usepackage{enumitem}
\usepackage{hyperref}
\usepackage{xcolor}
\usepackage{setspace}
\usepackage{fancyhdr}
\usepackage{listings}

\geometry{margin=2.5cm}

\titleformat{\section}{\Large\bfseries}{\thesection}{1em}{}
\titleformat{\subsection}{\bfseries}{\thesubsection}{1em}{}

\lstset{
    basicstyle=\ttfamily\small,
    backgroundcolor=\color{gray!10},
    frame=single,
    breaklines=true,
    showstringspaces=false
}

\begin{document}

\begin{center}
    \centering

    {\huge \textbf{Análise, Projeto e Desenvolvimento Ágil}}\\[0.5cm]

    {\Huge \textbf{Glossário Downstream}}\\[0.5cm]

    {\large Aluno 1: \textbf{Henrique Maia Cardosa}}\\[0.5cm]
    {\large Aluno 2: \textbf{João Miguel de Castro Menna}}\\[0.5cm]
    {\large Aluno 3: \textbf{João Pedro Izidoro}}\\[0.5cm]
    {\large Aluno 4: \textbf{Ricardo Gabriel Fialho Santos}}\\[0.5cm]
\end{center}

\vspace{1em}
\hrule
\vspace{1em}

\section*{Glossário}

\paragraph{}
\textbf{Cobertura de Testes:} É uma métrica utilizada na Engenharia de Software que avalia a eficácia dos testes realizados no código fonte. Ela mede em porcentagem a quantidade de código testado, assim é possível verificar o quão bem os testes estão se saindo.

\paragraph{}
\textbf{SAST (Static Application Security Testing):} O Teste de segurança de aplicação estática é um teste que verifica os padrões de  vulnerabilidades de segurança no código.

\paragraph{}
\textbf{Quality gates:} É a barreira de métricas necessárias para passar o código para frente.

\paragraph{}
\textbf{CI (Integração Contínua):} Integrar o código novo a todo momento na aplicação.

\paragraph{}
\textbf{CD (Entrega/Implantação Contínua):} Entregar o código para o cliente pronto para execução a cada certo período de tempo ou commit.

\paragraph{}
\textbf{Política de versionamento:} É um conjunto de regras estratégicas que visam gerenciar as mudanças de versão de uma aplicação, documento, código, ao longo do tempo.

\paragraph{}
\textbf{Dockerfile (práticas mínimas):} Arquivo que contém a especificação do contêiner para ser executado, pode possuir build de múltiplos estágios, variáveis de ambiente, ENTRYPOINT para aplicações CLI, CMD para aplicações que rodam apenas uma vez, por aí vai. É necessário adicionar o conteúdo ao container usando COPY.

\paragraph{}
\textbf{Docker-compose:} Aplicativo para a orquestração de containers Docker. Ele pode determinar quem subir primeiro, quem executa depois, com qual comando, com qual health check, etc.

\paragraph{}
\textbf{Estratégias de release:} blue-green, canário (canary), feature flags: São técnicas utilizadas no desenvolvimento de software e práticas de DevOps para minimizar riscos, reduzir o tempo de inatividade (downtime) e controlar a exposição de novas funcionalidades aos usuários.

\paragraph{}
\textbf{Plano de rollback / rollback:} É o processo de reverter um sistema, aplicação ou banco de dados para um estado anterior e funcional.

\paragraph{}
\textbf{Observabilidade:} Capacidade de entender o estado interno de um sistema a partir de suas saídas externas, como logs, métricas e traces.

\paragraph{}
\textbf{Logs estruturados:} Registros de eventos que seguem um formato predefinido e organizado, como JSON, diferentemente dos logs de texto simples

\paragraph{}
\textbf{Métricas (ex.: erros/min, tempo de resposta, uso de recursos):} São medidas quantitativas para monitorar, avaliar o desempenho de um sistema ou aplicação.

\paragraph{}
\textbf{Rastros/traços (tracing):} É acompanhar o decorrer de um projeto, desde observar, visualizar, localizar e depurar.

\paragraph{}
\textbf{Indicadores DORA (lead time, frequência de deploy, MTTR, taxa de falha):} São um conjunto de métricas-chave que avaliam o desempenho de equipes de desenvolvimento de software em práticas DevOps, visando a melhoria da eficiência e estabilidade. Frequência de implantação, tempo de lead para mudanças, taxa de falha e tempo médio de recuperação são as métricas.

\paragraph{}
\textbf{Runbook:} É um guia prático e detalhado com instruções passo a passo para executar tarefas operacionais, responder a incidentes ou manter sistemas. Serve como um manual de procedimentos, documentando o que fazer, como fazer e a ordem - para garantir que qualquer pessoa consiga executar o procedimento de forma consistente e segura.

\paragraph{}
\textbf{Post-mortem:} Prática que habilita a cultura do aprendizado com a falha e o blameless (sem culpa). O objetivo é entender o que aconteceu, identificar causas-raiz e aprender com a experiência.

\paragraph{}
\textbf{Governança da qualidade:} É a gestão da qualidade, que envolve diretrizes, transparência e conformidade nos processos, no contexto da engenharia de software isso se aplica a um projeto de uma aplicação.

\paragraph{}
\textbf{E2E (End-to-End):} Prática de testar a aplicação de ponta a ponta, passando por todas as partes do sistema.

\end{document}

